\section{Adição de símbolos por definição}\label{sec:adicaoSimbolosPorDefinicao}
\index{definição!extensão por definição}

A Teoria dos Conjuntos é uma teoria cuja linguagem não lógica inicial consiste apenas no símbolo de pertencimento ($\in$).
Conforme veremos, precisaremos rapidamente de símbolos funcionais adicionais, como o símbolo de constante $\emptyset$ (que representa o conjunto vazio), o operador binário de união de dois conjuntos ($\cup$) e o operador unário do conjunto das partes ($\mathcal{P}$).
Analogamente, também passaremos a utilizar símbolos relacionais adicionais, como o binário $\subseteq$.

Neste texto, tais símbolos serão introduzidos \emph{por definição}.
Intuitivamente, isto significa que eles servem apenas como abreviações: em vez de escrever repetidamente uma fórmula longa, introduzimos um símbolo novo cujo comportamento é \emph{explicitamente} fixado por um axioma.
O ponto central desta seção é formalizar essa prática e justificar que ela não aumenta o poder de prova da teoria em relação às sentenças da linguagem original.

\subsection{O quantificador \texorpdfstring{$\exists!$}{∃!}}\label{subsec:quantificadorExistsUnique}
\index{quantificador!existência única}

Os quantificadores primitivos da lógica de primeira ordem, neste texto, consistem apenas do quantificador universal ($\forall$) e do quantificador existencial ($\exists$).

Um quantificador amplamente usado em Matemática é o quantificador existencial único, denotado por $\exists!$.
Intuitivamente, $\exists! x \varphi$ significa que existe exatamente um $x$ tal que $\varphi$ é verdadeira.

Introduziremos este operador como uma abreviação, conforme explicitado na definição a seguir.

\begin{metadefinition}[Quantificador de existência única]\label{mdef:quantificadorExistenciaUnica}\index{existência única}
    Seja $\varphi$ uma fórmula.
    A notação $\exists!x\,\varphi$ é uma abreviação para a fórmula
    \begin{equation}\label{eq:existsUnique}
        (\exists x\,\varphi)\ \wedge\ \forall x\,\forall y\bigl((\varphi \wedge \varphi[x/y])\rightarrow x=y\bigr),
    \end{equation}
        em que $y$ é a primeira variável distinta de $x$ que não ocorre em $\varphi$.
\end{metadefinition}

Na Metadefinição~\ref{mdef:quantificadorExistenciaUnica}, a variável $y$ é tomada como uma variável auxiliar.
Não seria razoável que a escolha de tal variável dependesse da letra escolhida para representá-la.
Isso é formalizado pelo resultado a seguir.


\begin{metalemma}[Independência do auxiliar em $\exists!$]\label{mlem:independenciaAuxiliarExistsUnique}
    Seja $\varphi$ uma fórmula.
    Se $y$ e $z$ são variáveis distintas de $x$, distintas entre si, e que não ocorrem em $\varphi$, então as duas fórmulas obtidas ao expandir
    $\exists!x\,\varphi$ pela Equação~\eqref{eq:existsUnique} usando $y$ ou usando $z$ são equivalentes em qualquer teoria de primeira ordem.
\end{metalemma}

\begin{proof}
    Seja $T$ uma teoria de primeira ordem.
    Sejam $y, z$ variáveis que não ocorrem em $\varphi$, distintas de $x$ e distintas entre si.
    Suponha que $T\vdash (\exists x\,\varphi)\ \wedge\ \forall x\,\forall y\bigl((\varphi \wedge \varphi[x/y])\rightarrow x=y\bigr)$.
    Veremos que $T\vdash (\exists x\,\varphi)\ \wedge\ \forall x\,\forall z\bigl((\varphi \wedge \varphi[x/z])\rightarrow x=z\bigr)$.
    A recíproca é análoga.

    A parte existencial $\exists x\,\varphi$ já é a mesma nas duas expansões.
    Para a parte de unicidade, pela comutatividade dos quantificadores universais obtemos
    \[
        T\vdash \forall y\,\forall x\bigl((\varphi \wedge \varphi[x/y])\rightarrow x=y\bigr).
    \]
    Como $z$ não ocorre em $\varphi$, podemos renomear o quantificador externo de $y$ para $z$.
    Usando o Metalema~\ref{mlem:renomeacaoVariavelUniversal} e a igualdade
    $\varphi[x/y][y/z]=\varphi[x/z]$, concluímos
    \[
        T\vdash \forall z\,\forall x\bigl((\varphi \wedge \varphi[x/z])\rightarrow x=z\bigr).
    \]
    Aplicando novamente a comutatividade de quantificadores universais, segue
    \[
        T\vdash \forall x\,\forall z\bigl((\varphi \wedge \varphi[x/z])\rightarrow x=z\bigr),
    \]
    como queríamos.
\end{proof}

\subsection{Definindo definições}\label{subsec:definindoDefinicoes}
Agora estamos prontos para falar formalmente sobre definições.

\begin{metadefinition}[Definição explícita]\label{mdef:definicaoExplicita}\index{definição!explícita}\index{símbolo!não lógico!definido}
Seja $T$ uma teoria de primeira ordem cujo léxico é $\mathcal L$, e seja $\mathcal L'$ um léxico que estende $\mathcal L$.

\begin{enumerate}[label=(\roman*)]
    \item \textbf{Predicados.}
    Seja $P\in\mathcal L'\setminus\mathcal L$ um símbolo de predicado $n$-ário.
    Uma \emph{$T$-definição} de $P$ é uma sentença da forma
    \[
        \forall x_1\ldots \forall x_n\bigl(P(x_1,\ldots,x_n)\leftrightarrow \varphi(x_1,\ldots,x_n)\bigr),
    \]
    em que $x_1,\ldots,x_n$ são variáveis duas a duas distintas e $\varphi$ é uma $\mathcal L$-fórmula cujas variáveis livres estão entre $x_1,\ldots,x_n$.
    No caso $n=0$, uma $T$-definição de $P$ é uma sentença $P\leftrightarrow \varphi$,
    com $\varphi$ uma $\mathcal L$-sentença.

    \item \textbf{Funções.}
    Seja $f\in\mathcal L'\setminus\mathcal L$ um símbolo de função $n$-ário, com $n\geq 1$.
    Uma \emph{$T$-definição} de $f$ é uma sentença da forma
    \[
        \forall x_1\ldots \forall x_n\,\forall y\bigl(f(x_1,\ldots,x_n)=y \leftrightarrow \varphi(x_1,\ldots,x_n,y)\bigr),
    \]
    em que $x_1,\ldots,x_n,y$ são variáveis duas a duas distintas e $\varphi$ é uma $\mathcal L$-fórmula cujas variáveis livres estão entre $x_1,\ldots,x_n,y$, tal que
    \[
        T\vdash \forall x_1\ldots \forall x_n\,\exists!y\,\varphi(x_1,\ldots,x_n,y).
    \]
    (Intuitivamente, $T$ prova que $\varphi$ define \emph{unicamente} um valor $y$ para cada entrada $x_1,\ldots,x_n$.)

    \item \textbf{Constantes.}
    Seja $c\in\mathcal L'\setminus\mathcal L$ um símbolo de constante.
    Uma \emph{$T$-definição} de $c$ é uma sentença da forma
    \[
        \forall y\bigl(c=y \leftrightarrow \varphi(y)\bigr),
    \]
    em que $\varphi$ é uma $\mathcal L$-fórmula cujas variáveis livres estão entre $y$, tal que $T\vdash \exists!y\,\varphi(y)$.
\end{enumerate}

Em cada um dos casos acima, dizemos que $\varphi$ é a \emph{fórmula definidora} do símbolo novo.
Quando for conveniente, escreveremos essa fórmula como
\[
\Def_P(x_1,\ldots,x_n),\qquad \Def_f(x_1,\ldots,x_n,y),\qquad \Def_c(y),
\]
conforme o caso.

Além disso, se $P$, $f$ ou $c$ são símbolos dos tipos correspondentes em $\mathcal L$, convencionamos:

\begin{itemize}
    \item $\Def_P(x_1,\ldots,x_n)$ é a fórmula $P(x_1,\ldots,x_n)$.
    No caso $n=0$, $\Def_P$ é a fórmula $P$.
    \item $\Def_f(x_1,\ldots,x_n,y)$ é a fórmula $f(x_1,\ldots,x_n)=y$.
    \item $\Def_c(y)$ é a fórmula $c=y$.
    \item Se $z$ é uma variável, então $\Def_{z}(z, y)$ é a fórmula $y=z$.
\end{itemize}

Note que, nesses casos, $\Def_P$, $\Def_f$, $\Def_c$ e $\Def_{z}$ são fórmulas da linguagem de $\mathcal L$.
\end{metadefinition}

\begin{example}\label{exa:definicoesExplicitas}
Ainda que não tenhamos apresentado a Teoria dos Conjuntos, vejamos alguns exemplos.
\begin{itemize}
    \item \textbf{Teoria dos conjuntos.}
    Uma definição natural de $x\subseteq y$ na linguagem de $\in$ é:
    \[
        \forall x\,\forall y\Bigl(x\subseteq y\ \leftrightarrow\ \forall z\,(z\in x\rightarrow z\in y)\Bigr).
    \]
    Aqui $\subseteq$ é um novo símbolo de predicado binário, definido por uma $\mathcal L$-fórmula (em que $\mathcal L=\{\in\}$).
    Nesse caso, $\Def_{\subseteq}$ é a fórmula $\forall z\,(z\in x\rightarrow z\in y)$.

    \item \textbf{Teoria dos grupos.}
    Considere a teoria de grupos na linguagem \emph{sem} símbolo de constante para a identidade,
    isto é, apenas com o símbolo de função binária $\cdot$.
    Uma definição explícita (quando possível em uma teoria dada) para um símbolo novo de constante $e$ é:
    \[
        \forall y\Bigl(e=y\ \leftrightarrow\ \forall x\,\bigl(x\cdot y=x \wedge y\cdot x=x\bigr)\Bigr).
    \]
    A condição de aplicabilidade é que a teoria em questão prove $\exists!y\,\forall x(x\cdot y=x\wedge y\cdot x=x)$.
    Em uma axiomatização usual de grupos \emph{sem} constante, isso é demonstrável.
    Nesse caso, $\Def_e$ é a fórmula $\forall x\,\bigl(x\cdot y=x \wedge y\cdot x=x\bigr)$.
\end{itemize}
\end{example}

Agora podemos definir a noção de extensão por definição, uma formalização da prática de introduzir definições como abreviações.

\begin{metadefinition}[Extensão por definição]\label{mdef:extensaoPorDefinicao}\index{definição!extensão por definição}
    Seja $T$ uma teoria de primeira ordem cujo léxico é $\mathcal L$.
    Dizemos que uma teoria $T'$ é uma \emph{extensão por definição} de $T$ se:
    \begin{itemize}
        \item o léxico de $T'$ é um léxico $\mathcal L'$ que estende $\mathcal L$;
        \item os axiomas de $T'$ consistem nos axiomas de $T$ e, para cada símbolo novo de $\mathcal L'\setminus\mathcal L$,
        em uma (única) definição explícita no sentido da Metadefinição~\ref{mdef:definicaoExplicita}.
    \end{itemize}
\end{metadefinition}

Essa noção pode causar estranhamento: estamos formalizando a introdução de definições como a adição de novos axiomas e, portanto, alterando a teoria base.
Isso é necessário no formalismo adotado, pois ele não captura diretamente a prática informal de introduzir definições como meras abreviações.

A seguir, veremos que extensões por definição não aumentam nem o poder de prova sobre sentenças da linguagem original, nem o poder de expressão.
Tudo o que é expressável ou demonstrável na linguagem estendida também é expressável ou demonstrável na linguagem original, respectivamente.
Portanto, apesar de, formalmente, estarmos adicionando novos símbolos e axiomas, a extensão por definição captura as propriedades essenciais da prática de introduzir definições como abreviações.


\subsection{Eliminação de símbolos definidos: termos}\label{subsec:eliminacaoSimbolosDefinidosTermos}
\index{eliminação de símbolos definidos}

A ideia central é associar a cada termo $t$ da linguagem $\mathcal L'$ uma fórmula da linguagem $\mathcal L$ que semanticamente descreve $t$.
Isto substitui termos com símbolos definidos por uma descrição na linguagem original.

\begin{metadefinition}[Fórmula do valor de um termo]\label{mdef:formulaDoValorDeUmTermo}\index{termo!fórmula do valor}\index{eliminação de símbolos definidos!termos}
    Suponha que $T'$ é uma extensão por definição de $T$, com léxicos $\mathcal L'$ e $\mathcal L$ respectivamente.
    Seja $t$ um termo de $\mathcal L'$.
    Definimos por recursão uma $\mathcal L$-fórmula $\delta_t$ cujas variáveis livres estão entre as variáveis de $t$ e mais uma variável $y$, lida como
    ``$y$ é o valor do termo $t$'', como segue.

    \begin{enumerate}[label=(\roman*)]
        \item Se $t$ é um símbolo de constante $c$, então $\delta_t(y)$ é a fórmula $\Def_c(y)$.
        \item Se $t$ é uma variável $z$, então $\delta_t(z,y)$ é a fórmula $y=z$, ou seja, $\Def_z(z,y)$.
        \item Se $t$ é $f(t_1,\ldots,t_m)$, então $\delta_t(x_1, \dots, x_n,y)$ é:
        \[
            \exists y_1\ldots \exists y_m\Bigl(
            \bigwedge_{i=1}^m \delta_{t_i}(x_1, \dots, x_n,y_i)\ \wedge\ \Def_f(y_1,\ldots,y_m,y)\Bigr),
        \]
        em que $x_1, \dots, x_n$ é uma lista que contém todas as variáveis que ocorrem em $t$,
        e $y_1, \dots, y_m$ são variáveis auxiliares duas a duas distintas e distintas de $x_1,\ldots,x_n,y$.
    \end{enumerate}
\end{metadefinition}

\begin{metalemma}[Eliminação de termos]\label{mlem:eliminacaoDeTermos}\index{eliminação de símbolos definidos!termos}
    Seja $T'$ uma extensão por definição de $T$, com $\mathcal L\subseteq \mathcal L'$.
    Seja $t$ um termo de $\mathcal L'$ cujas variáveis estejam entre $x_1,\ldots,x_n$,
    e seja $y$ uma variável distinta de $x_1,\ldots,x_n$.
    Seja $\delta_t(x_1,\ldots,x_n,y)$ como na Metadefinição~\ref{mdef:formulaDoValorDeUmTermo}. Então:
    \begin{enumerate}[label=(\arabic*)]
        \item $T'\vdash \forall x_1\ldots \forall x_n\,\forall y\bigl(t=y \leftrightarrow \delta_t(x_1,\ldots,x_n,y)\bigr)$.
        \item $T\vdash \forall x_1\ldots \forall x_n\,\exists!y\,\delta_t(x_1,\ldots,x_n,y)$.
    \end{enumerate}

    Além disso, se $t$ é um $\mathcal L$-termo, então:
    \begin{enumerate}[label=(\arabic*)]
        \setcounter{enumi}{2}
        \item $T\vdash \forall x_1\ldots \forall x_n\,\forall y\bigl(t=y \leftrightarrow \delta_t(x_1,\ldots,x_n,y)\bigr)$.
    \end{enumerate}
\end{metalemma}

\begin{proof}
Provaremos (1), (2) e (3) simultaneamente por indução sobre a complexidade do termo $t$.
Escreva $\bar x$ para $x_1,\ldots,x_n$.

\smallskip
\textbf{Base 1: $t$ é uma constante $c$.}
 Pela Metadefinição~\ref{mdef:formulaDoValorDeUmTermo}, $\delta_t(y)$ é $\Def_c(y)$.

\emph{(1).} Se $c\in\mathcal L'\setminus\mathcal L$, então a sentença definidora de $c$ pertence a $T'$ e tem a forma
$\forall y\,(c=y\leftrightarrow \Def_c(y))$, que é exatamente (1).
Se $c\in\mathcal L$, então por convenção $\Def_c(y)$ é $c=y$, logo
$\forall y\,(c=y\leftrightarrow \Def_c(y))$ é uma instância de tautologia.

\emph{(2).} Se $c$ é novo, a própria hipótese de ``introdução por definição'' garante
$T\vdash \exists!y\,\Def_c(y)$.
Se $c\in\mathcal L$, então $\Def_c(y)$ é $c=y$; provamos $T\vdash \exists!y\,(c=y)$:
\begin{itemize}
    \item \emph{Existência.} Pela reflexividade da igualdade, $T\vdash c=c$.
    Aplicando FO15 à fórmula $\phi(y)\equiv (c=y)$ com o termo $t=c$, obtemos o axioma lógico
    $(c=c)\rightarrow \exists y(c=y)$; por \emph{modus ponens}, $T\vdash \exists y(c=y)$.
    \item \emph{Unicidade.} Dos axiomas de igualdade, temos $c=y\rightarrow y=c$ (simetria) e
    $y=c\wedge c=z\rightarrow y=z$ (transitividade). Assim,
    $(c=y\wedge c=z)\rightarrow y=z$, e portanto a parte de unicidade de $\exists!y(c=y)$.
\end{itemize}

\emph{(3).} Se $t$ é um $\mathcal L$-termo, então necessariamente $c\in\mathcal L$ e, como acima,
$\forall y\,(c=y\leftrightarrow \Def_c(y))$ é tautológica; logo vale em $T$.

\smallskip
\textbf{Base 2: $t$ é uma variável $z$.}
Pela Metadefinição~\ref{mdef:formulaDoValorDeUmTermo}, $\delta_t(z,y)$ é $y=z$, isto é, $\Def_{z}(z,y)$.

\emph{(1).} Queremos $T'\vdash \forall z\,\forall y\,(z=y\leftrightarrow y=z)$.
Pela simetria da igualdade, temos
$T'\vdash z=y\rightarrow y=z$ e também $T'\vdash y=z\rightarrow z=y$.
Por uma instância da tautologia proposicional
\[
(P\rightarrow Q)\rightarrow\bigl((Q\rightarrow P)\rightarrow (P\leftrightarrow Q)\bigr),
\]
obtemos $T'\vdash z=y\leftrightarrow y=z$, e tomando fecho universal em $z,y$ concluímos (1).

\emph{(2).} Provamos $T\vdash \forall z\,\exists!y\,(y=z)$.
A existência segue de $T\vdash z=z$ e de FO15 aplicado a $\phi(y)\equiv (y=z)$ com o termo $t=z$.
A unicidade segue de $(y=z\wedge w=z)\rightarrow y=w$, obtida por simetria e transitividade da igualdade.
Logo vale (2) no caso variável.

\emph{(3).} Como variáveis já pertencem à linguagem, o argumento de (1) funciona igualmente dentro de $T$.

\smallskip
\textbf{Passo indutivo: $t=f(t_1,\ldots,t_m)$.}
Sejam $\bar x$ as variáveis que ocorrem em $t$ (no sentido da sua convenção: todas as variáveis livres estão dentre $\bar x$).
Pela Metadefinição~\ref{mdef:formulaDoValorDeUmTermo},
\[
\delta_t(\bar x,y)\ \equiv\ \exists y_1\ldots\exists y_m\Bigl(\bigwedge_{i=1}^m \delta_{t_i}(\bar x,y_i)\ \wedge\ \Def_f(y_1,\ldots,y_m,y)\Bigr).
\]

\smallskip
\emph{(1) Trabalhando em $T'$.}
Pela hipótese de indução aplicada a cada $t_i$, temos
\[
T'\vdash \forall \bar x\,\forall y_i\bigl(t_i=y_i\leftrightarrow \delta_{t_i}(\bar x,y_i)\bigr)\qquad (i=1,\ldots,m).
\]
Usando essas equivalências dentro do existencial (raciocínio proposicional), obtemos em $T'$:
\[
\delta_t(\bar x,y)\ \leftrightarrow\
\exists \bar y\Bigl(\bigwedge_{i=1}^m (t_i=y_i)\ \wedge\ \Def_f(\bar y,y)\Bigr).
\]
Agora, em $T'$ temos sempre
\[
\forall \bar u\,\forall v\bigl(f(\bar u)=v\leftrightarrow \Def_f(\bar u,v)\bigr),
\]
pois ou $f\in\mathcal L$ e então isto é tautológico pela convenção $\Def_f(\bar u,v)\equiv f(\bar u)=v$,
ou $f$ é novo e esta é precisamente a sentença definidora de $f$.
Logo podemos substituir $\Def_f(\bar y,y)$ por $f(\bar y)=y$ dentro de $T'$.

Com isso, e usando a congruência para funções e a transitividade da igualdade,
temos a equivalência (em $T'$):
\[
f(t_1,\ldots,t_m)=y\ \leftrightarrow\ \exists \bar y\Bigl(\bigwedge_{i=1}^m (t_i=y_i)\wedge f(\bar y)=y\Bigr).
\]
Juntando com a equivalência anterior para $\delta_t(\bar x,y)$, concluímos
\[
T'\vdash \forall \bar x\,\forall y\bigl(t=y\leftrightarrow \delta_t(\bar x,y)\bigr),
\]
isto é, (1).

\smallskip
\emph{(2) Trabalhando em $T$.}
Pela hipótese de indução, para cada $i\le m$,
\[
T\vdash \forall\bar x\,\exists!y_i\,\delta_{t_i}(\bar x,y_i).
\]
Fixe $\bar x$. Pela parte existencial, existem (únicos) $y_1,\ldots,y_m$ tais que $\delta_{t_i}(\bar x,y_i)$.
Agora mostramos que existe um único $y$ tal que $\Def_f(y_1,\ldots,y_m,y)$:
\begin{itemize}
    \item Se $f$ é novo, isto é exatamente a condição da Metadefinição~\ref{mdef:definicaoExplicita}:
    $T\vdash \forall \bar u\,\exists!v\,\Def_f(\bar u,v)$.
    \item Se $f\in\mathcal L$, então $\Def_f(\bar u,v)$ é $f(\bar u)=v$.
    Para a existência, a reflexividade da igualdade dá $T\vdash f(\bar u)=f(\bar u)$,
    e então FO15, aplicado à fórmula $\phi(v)\equiv (f(\bar u)=v)$ com o termo $t=f(\bar u)$, dá $\exists v\,f(\bar u)=v$.
    Para a unicidade, de $f(\bar u)=v$ e $f(\bar u)=w$ obtemos $v=w$ por simetria e transitividade da igualdade.
\end{itemize}
Portanto, existe um único $y$ tal que $\Def_f(y_1,\ldots,y_m,y)$.
Logo existe um único $y$ tal que
\[
\exists y_1\ldots\exists y_m\Bigl(\bigwedge_{i=1}^m \delta_{t_i}(\bar x,y_i)\ \wedge\ \Def_f(y_1,\ldots,y_m,y)\Bigr),
\]
isto é, $\exists!y\,\delta_t(\bar x,y)$.
Como $\bar x$ foi arbitrário, concluímos (2):
\[
T\vdash \forall \bar x\,\exists!y\,\delta_t(\bar x,y).
\]

\smallskip
\emph{(3) Caso adicional: $t$ é um $\mathcal L$-termo.}
Se $t$ é um $\mathcal L$-termo, então $f\in\mathcal L$ e cada $t_i$ também é $\mathcal L$-termo.
Pela hipótese de indução (3) aplicada a cada $t_i$, temos as equivalências $t_i=y_i\leftrightarrow \delta_{t_i}(\bar x,y_i)$ agora em $T$.
Repetindo o argumento de (1) inteiramente dentro de $T$ (e usando que $\Def_f(\bar y,y)$ é $f(\bar y)=y$ por convenção),
obtemos
\[
T\vdash \forall \bar x\,\forall y\bigl(t=y\leftrightarrow \delta_t(\bar x,y)\bigr),
\]
isto é, (3).

Isso conclui a indução e a prova do metalema.
\end{proof}

\subsection{Eliminação de símbolos definidos: fórmulas}\label{subsec:eliminacaoSimbolosDefinidosFormulas}
\index{eliminação de símbolos definidos}

Agora eliminamos símbolos definidos de fórmulas.

\begin{metadefinition}[Tradução $\psi\mapsto \psi^\ast$]\label{mdef:traducaoFormulas}\index{eliminação de símbolos definidos!fórmulas}
    Seja $\psi$ uma fórmula da linguagem $\mathcal L'$ cujas variáveis livres estejam entre $x_1,\ldots,x_n$.
    Definimos por recursão uma $\mathcal L$-fórmula $\psi^\ast$ (com as mesmas variáveis livres) como segue.
    Em todos os casos atômicos abaixo, as variáveis auxiliares $u,v,u_1,\ldots,u_m$ são escolhidas novas e distintas das variáveis livres listadas.

    \begin{enumerate}[label=(\roman*)]
        \item Se $\psi$ é $t_1=t_2$, então
        \[
            \psi^\ast\ :=\ \exists u\,\exists v\bigl(\delta_{t_1}(\bar x,u)\wedge \delta_{t_2}(\bar x,v)\wedge u=v\bigr).
        \]

        \item Se $\psi$ é um predicado $0$-ário $R\in\mathcal L$, então $\psi^\ast:=R$.

        \item Se $\psi$ é $R(t_1,\ldots,t_m)$, em que $R\in\mathcal L$ é predicado $m$-ário e $m>0$, então
        \[
            \psi^\ast\ :=\ \exists u_1\ldots \exists u_m\Bigl(\bigwedge_{i=1}^m \delta_{t_i}(\bar x,u_i)\ \wedge\ R(u_1,\ldots,u_m)\Bigr).
        \]

        \item Se $\psi$ é um predicado $0$-ário $P\in\mathcal L'\setminus\mathcal L$ definido por $\Def_P$, então $\psi^\ast:=\Def_P$.

        \item Se $\psi$ é $P(t_1,\ldots,t_m)$, em que $P\in\mathcal L'\setminus\mathcal L$ é predicado $m$-ário definido por $\Def_P$ e $m>0$, então
        \[
            \psi^\ast\ :=\ \exists u_1\ldots \exists u_m\Bigl(\bigwedge_{i=1}^m \delta_{t_i}(\bar x,u_i)\ \wedge\ \Def_P(u_1,\ldots,u_m)\Bigr).
        \]

        \item Conectivos e quantificadores:
        \[
            (\neg\theta)^\ast:=\neg(\theta^\ast),\quad
            (\theta\square\eta)^\ast:=\theta^\ast\square\eta^\ast,
        \]
        para $\square\in\{\wedge,\vee,\rightarrow,\leftrightarrow\}$, e
        \[
            (\forall z\,\theta)^\ast:=\forall z\,(\theta^\ast),\qquad
            (\exists z\,\theta)^\ast:=\exists z\,(\theta^\ast).
        \]
    \end{enumerate}
\end{metadefinition}

\begin{metalemma}[Eliminação de fórmulas]\label{mlem:eliminacaoDeFormulas}\index{eliminação de símbolos definidos!fórmulas}
    Seja $T'$ uma extensão por definição de $T$.
    Para toda fórmula $\psi(x_1, \dots, x_n)$ da linguagem $\mathcal L'$,
    existe uma fórmula $\psi^\ast$ na linguagem $\mathcal L$ (como na Metadefinição~\ref{mdef:traducaoFormulas})
    com as mesmas variáveis livres que $\psi$ tal que
    \[
        T'\vdash \forall x_1\ldots \forall x_n\,(\psi\leftrightarrow \psi^\ast).
    \]

    Além disso, se $\psi$ é uma fórmula da linguagem $\mathcal L$, então:
    \[
        T\vdash \forall x_1\ldots \forall x_n\,(\psi\leftrightarrow \psi^\ast).
    \]
\end{metalemma}

\begin{proof}
    Provaremos por indução sobre a estrutura de $\psi$ a seguinte afirmação simultânea.

    \smallskip
    \textbf{Afirmação A.} Para toda fórmula $\psi(\bar x)$ de $\mathcal L'$, vale
    \[
        T'\vdash \forall \bar x\,(\psi\leftrightarrow \psi^\ast).
    \]

    \textbf{Afirmação B.} Se, além disso, $\psi(\bar x)$ é uma fórmula de $\mathcal L$, então
    \[
        T\vdash \forall \bar x\,(\psi\leftrightarrow \psi^\ast).
    \]

    \smallskip
    Em cada caso, $\psi^\ast$ é a tradução definida na Metadefinição~\ref{mdef:traducaoFormulas}.

    \smallskip
    \textbf{Caso atômico 1:} $\psi$ é $t_1=t_2$.

    \emph{Prova de A.}
    Pelo item (1) do Metalema~\ref{mlem:eliminacaoDeTermos},
    \[
        T'\vdash \forall \bar x\,\forall u\,(t_1=u\leftrightarrow \delta_{t_1}(\bar x,u))
        \quad\text{e}\quad
        T'\vdash \forall \bar x\,\forall v\,(t_2=v\leftrightarrow \delta_{t_2}(\bar x,v)).
    \]
    Usando raciocínio proposicional (instâncias de tautologias) e substitutividade da igualdade, obtemos em $T'$:
    \[
        T'\vdash \forall \bar x\Bigl(t_1=t_2 \ \leftrightarrow\ \exists u\,\exists v\bigl(\delta_{t_1}(\bar x,u)\wedge \delta_{t_2}(\bar x,v)\wedge u=v\bigr)\Bigr),
    \]
    e o lado direito é, por definição, $\psi^\ast$.

    \emph{Prova de B.}
    Suponha agora que $\psi$ é uma $\mathcal L$-fórmula; então $t_1$ e $t_2$ são $\mathcal L$-termos.
    Pelo item (3) do Metalema~\ref{mlem:eliminacaoDeTermos},
    \[
        T\vdash \forall \bar x\,\forall u\,(t_1=u\leftrightarrow \delta_{t_1}(\bar x,u))
        \quad\text{e}\quad
        T\vdash \forall \bar x\,\forall v\,(t_2=v\leftrightarrow \delta_{t_2}(\bar x,v)).
    \]
    O mesmo raciocínio proposicional usado acima (agora dentro de $T$) dá
    \[
        T\vdash \forall \bar x\,(\psi\leftrightarrow \psi^\ast).
    \]

    \smallskip
    \textbf{Caso atômico 2:} $\psi$ é uma fórmula atômica com símbolo de predicado $R\in\mathcal L$.

    Se $R$ tem aridade $0$, então $\psi^\ast$ é a própria fórmula $R$, e A e B seguem por tautologia e generalização.
    Suponha, portanto, que $\psi$ é $R(t_1,\ldots,t_m)$ com $m>0$.

    \emph{Prova de A.}
    Pelo item (1) do Metalema~\ref{mlem:eliminacaoDeTermos}, para cada $i\le m$,
    \[
        T'\vdash \forall \bar x\,\forall u_i\,(t_i=u_i\leftrightarrow \delta_{t_i}(\bar x,u_i)).
    \]
    Usando essas equivalências e a congruência para predicados,
    obtemos em $T'$:
    \[
        T'\vdash \forall \bar x\Bigl(R(t_1,\ldots,t_m)\ \leftrightarrow\
        \exists u_1\ldots \exists u_m\bigl(\bigwedge_{i=1}^m \delta_{t_i}(\bar x,u_i)\wedge R(u_1,\ldots,u_m)\bigr)\Bigr),
    \]
    e o lado direito é exatamente $\psi^\ast$.

    \emph{Prova de B.}
    Se $\psi$ é uma $\mathcal L$-fórmula, então cada $t_i$ é um $\mathcal L$-termo.
    Pelo item (3) do Metalema~\ref{mlem:eliminacaoDeTermos}, para cada $i$,
    \[
        T\vdash \forall \bar x\,\forall u_i\,(t_i=u_i\leftrightarrow \delta_{t_i}(\bar x,u_i)).
    \]
    Repetindo o mesmo argumento (congruência para $R$ e raciocínio proposicional), agora em $T$,
    concluímos
    \[
        T\vdash \forall \bar x\,(\psi\leftrightarrow \psi^\ast).
    \]

    \smallskip
    \textbf{Caso atômico 3:} $\psi$ é uma fórmula atômica com símbolo de predicado $P\in\mathcal L'\setminus\mathcal L$.

    Aqui somente A é relevante (B não se aplica, pois $\psi$ não está em $\mathcal L$).
    Se $P$ tem aridade $0$, então $\psi$ é $P$, $\psi^\ast$ é $\Def_P$, e a conclusão segue do axioma definidor $P\leftrightarrow\Def_P$ e de generalização.
    Suponha, portanto, que $\psi$ é $P(t_1,\ldots,t_m)$ com $m>0$.

    Pelo axioma definidor de $P$ em $T'$, temos
    \[
        T'\vdash \forall \bar u\,(P(\bar u)\leftrightarrow \Def_P(\bar u)).
    \]
    Combinando isso com o item (1) do Metalema~\ref{mlem:eliminacaoDeTermos} aplicado a cada $t_i$
    e com raciocínio proposicional, obtemos
    \[
        T'\vdash \forall \bar x\Bigl(P(t_1,\ldots,t_m)\leftrightarrow
        \exists u_1\ldots \exists u_m\bigl(\bigwedge_{i=1}^m \delta_{t_i}(\bar x,u_i)\wedge \Def_P(u_1,\ldots,u_m)\bigr)\Bigr),
    \]
    que é precisamente $\forall\bar x(\psi\leftrightarrow \psi^\ast)$.

    \smallskip
    \textbf{Conectivos.}
    Suponha que A (e, quando aplicável, B) valem para $\theta$ e $\eta$.

    Para a negação, da hipótese indutiva $T'\vdash \forall\bar x(\theta\leftrightarrow \theta^\ast)$
    e de uma instância de tautologia proposicional
    $(A\leftrightarrow B)\rightarrow(\neg A\leftrightarrow \neg B)$,
    obtemos $T'\vdash \forall\bar x(\neg\theta\leftrightarrow \neg\theta^\ast)$,
    e pela definição $(\neg\theta)^\ast=\neg(\theta^\ast)$, isto é exatamente A para $\neg\theta$.
    O mesmo argumento vale em $T$ quando $\theta$ está em $\mathcal L$, provando B.

    Para os conectivos binários $\square\in\{\wedge,\vee,\rightarrow,\leftrightarrow\}$,
    usamos instâncias de tautologias do tipo
    \[
        (A\leftrightarrow A')\wedge(B\leftrightarrow B')\ \rightarrow\ ((A\square B)\leftrightarrow(A'\square B')),
    \]
    junto com as hipóteses indutivas para $\theta$ e $\eta$,
    e a definição $(\theta\square\eta)^\ast=\theta^\ast\square\eta^\ast$.
    Isso prova A (e também B quando $\theta,\eta$ são fórmulas de $\mathcal L$).

    \smallskip
    \textbf{Quantificadores.}
    Suponha que A vale para $\theta$.
    No caso de $\forall z\,\theta$, aplicamos a hipótese indutiva a $\theta$ considerando suas variáveis livres entre $\bar x,z$.
    Assim,
    \[
        T'\vdash \forall \bar x\,\forall z\,(\theta\leftrightarrow \theta^\ast).
    \]
    Pelas regras derivadas para quantificadores e para bicondicional, obtemos
    \[
        T'\vdash \forall \bar x\,(\forall z\,\theta \leftrightarrow \forall z\,\theta^\ast).
    \]
    Como $(\forall z\,\theta)^\ast=\forall z\,\theta^\ast$, isso prova A para $\forall z\,\theta$.

    O caso de $\exists z\,\theta$ é análogo, usando a equivalência
    \[
        \forall z\,(\theta\leftrightarrow\theta^\ast)
    \]
    para obter
    \[
        \exists z\,\theta\leftrightarrow \exists z\,\theta^\ast.
    \]

    O mesmo argumento, realizado em $T$, prova B quando a fórmula quantificada pertence à linguagem $\mathcal L$.

    \smallskip
    Isso conclui a indução e prova simultaneamente as duas afirmações, logo o metalema.
\end{proof}

\subsection{Extensões por definição são conservativas}\label{subsec:extensoesPorDefinicaoSaoConservativas}

\begin{metadefinition}[Extensão conservativa]\label{mdef:extensaoConservativa}\index{extensão!conservativa}
    Seja $T$ uma teoria de primeira ordem com léxico $\mathcal L$, e seja $T'$ uma teoria com léxico $\mathcal L'$ estendendo $\mathcal L$.
    Dizemos que $T'$ é uma \emph{extensão conservativa} de $T$ se cada axioma de $T$ é também um axioma de $T'$ e, para cada $\mathcal L$-sentença $\phi$,
    \[
        T'\vdash \phi \ \Longrightarrow\ T\vdash \phi.
    \]
\end{metadefinition}

O teorema a seguir formaliza a ideia de que definições são apenas abreviações.

\begin{metatheorem}\label{mthm:conservatividadeDaExtensaoPorDefinicao}\index{definição!extensão por definição!conservatividade}
    Seja $T$ uma teoria de primeira ordem e $T'$ uma extensão por definição de $T$.
    Então $T'$ é uma extensão conservativa de $T$.
\end{metatheorem}

\begin{proof}
    Seja $\phi$ uma $\mathcal L$-sentença e suponha $T'\vdash \phi$.
    Considere uma demonstração formal de $\phi$ a partir de $T'$,
    \[
        \theta_0,\theta_1,\ldots,\theta_n,
    \]
    em que $\theta_n$ é $\phi$.

    Para cada $i$, considere a fórmula $\theta_i^\ast$ obtida pela tradução da Metadefinição~\ref{mdef:traducaoFormulas}.
    Cada $\theta_i^\ast$ é uma fórmula da linguagem $\mathcal L$.

    \smallskip
    \textbf{Afirmação.} Para todo $i\le n$, vale $T\vdash \theta_i^\ast$.

    \smallskip
    Provaremos por indução em $i$.
    Suponha provado para todos $j<i$.

    \begin{itemize}
        \item Se $\theta_i$ é um axioma de $T$, então $\theta_i$ é uma $\mathcal L$-sentença.
        Pelo Metalema~\ref{mlem:eliminacaoDeFormulas} (aplicado a uma fórmula já na linguagem $\mathcal L$),
        temos $T\vdash \theta_i\leftrightarrow \theta_i^\ast$.
        Em particular, $\theta_i^\ast$ é consequência lógica proposicional de $\theta_i$ e de uma instância de tautologia;
        como $\theta_i$ é axioma de $T$, concluímos $T\vdash \theta_i^\ast$.

        \item Se $\theta_i$ é uma sentença definidora (axioma de definição) de algum símbolo novo, então,
        $T\vdash \theta_i^\ast$.
        De fato, a tradução substitui o símbolo definido pela sua fórmula definidora e substitui os termos variáveis por suas fórmulas de valor, que são equivalentes a igualdades triviais pelo Metalema~\ref{mlem:eliminacaoDeTermos}.
        Como as fórmulas definidoras já são fórmulas de $\mathcal L$, a segunda parte do Metalema~\ref{mlem:eliminacaoDeFormulas} identifica cada fórmula definidora com sua própria tradução dentro de $T$.
        Assim, no caso de predicados, a tradução da sentença definidora se reduz em $T$ a uma instância de $\alpha\leftrightarrow\alpha$; nos casos de funções e constantes, reduz-se do mesmo modo após eliminar o termo definido pelo lado que usa $\Def_f$ ou $\Def_c$.

        \item Se $\theta_i$ é um axioma lógico (no sentido do sistema fixado anteriormente),
        então $T\vdash \theta_i^\ast$.
        A verificação é sistemática: usa-se que $\ast$ comuta com conectivos e quantificadores e que termos são eliminados pelo Metalema~\ref{mlem:eliminacaoDeTermos}.
        Nos axiomas de congruência, se o símbolo envolvido é novo, usa-se ainda que a fórmula definidora é uma $\mathcal L$-fórmula, a Lei da igualdade de Leibniz e, no caso de símbolos funcionais, a hipótese de unicidade exigida na Metadefinição~\ref{mdef:definicaoExplicita}.
        Assim, $T\vdash \theta_i^\ast$.

        \item Se $\theta_i$ é obtida por \emph{modus ponens} a partir de $\theta_j$ e $\theta_k=\theta_j\rightarrow \theta_i$,
        com $j,k<i$, então por hipótese de indução $T\vdash \theta_j^\ast$ e $T\vdash (\theta_j\rightarrow \theta_i)^\ast$.
        Pela definição de $\ast$, $(\theta_j\rightarrow \theta_i)^\ast$ é $\theta_j^\ast\rightarrow \theta_i^\ast$.
        Logo, por \emph{modus ponens}, $T\vdash \theta_i^\ast$.

        \item Se $\theta_i$ é obtida pela regra de introdução do quantificador universal, então existem fórmulas $\chi,\eta$, uma variável $x$ e $j<i$ tais que
        $\theta_j$ é $\chi\rightarrow\eta$, $\theta_i$ é $\chi\rightarrow\forall x\,\eta$, e $x$ não ocorre livre em $\chi$.
        Pela hipótese de indução,
        \[
            T\vdash (\chi\rightarrow\eta)^\ast,
        \]
        isto é, $T\vdash \chi^\ast\rightarrow\eta^\ast$.
        Como $\chi^\ast$ tem as mesmas variáveis livres que $\chi$, a variável $x$ não ocorre livre em $\chi^\ast$.
        Pela regra de introdução do quantificador universal,
        \[
            T\vdash \chi^\ast\rightarrow\forall x\,\eta^\ast,
        \]
        que é exatamente $\theta_i^\ast$.

        \item Se $\theta_i$ é obtida pela regra de introdução do quantificador existencial, então existem fórmulas $\chi,\eta$, uma variável $x$ e $j<i$ tais que
        $\theta_j$ é $\chi\rightarrow\eta$, $\theta_i$ é $\exists x\,\chi\rightarrow\eta$, e $x$ não ocorre livre em $\eta$.
        Pela hipótese de indução, $T\vdash \chi^\ast\rightarrow\eta^\ast$.
        Como $\eta^\ast$ tem as mesmas variáveis livres que $\eta$, a variável $x$ não ocorre livre em $\eta^\ast$.
        Pela regra de introdução do quantificador existencial,
        \[
            T\vdash \exists x\,\chi^\ast\rightarrow\eta^\ast,
        \]
        que é exatamente $\theta_i^\ast$.
    \end{itemize}

    Isso conclui a indução, e em particular $T\vdash \theta_n^\ast=\phi^\ast$.

    \smallskip
    Resta observar que, como $\phi$ já está na linguagem $\mathcal L$ e é sentença,
    o Metalema~\ref{mlem:eliminacaoDeFormulas} dá $T\vdash \phi\leftrightarrow \phi^\ast$,
    e portanto $\phi$ é consequência proposicional de $\phi^\ast$.
    Como $T\vdash \phi^\ast$, segue $T\vdash \phi$.
\end{proof}
